\subsection{Triangularization over Complex and Real Fields}

\begin{definition} \label{def:diagonal-and-upper-triangular-matrices} \hlt{Diagonal and Upper-Triangular Matrices}\\
Let $n\in\mathbb{Z}_{>0}$ and $A=[a_{ij}]_{i,j=1}^{n}\in\operatorname{Mat}(n,n,\mathbb{F})$.
\begin{enumerate}[label=(\roman*)]
\setlength{\itemsep}{0pt}
\item The \hlt{diagonal} of $A$ is $(a_{11},\ldots,a_{nn})$.
\item The matrix $A$ is \hlt{diagonal} iff $(\forall\,i,j\in\{1,\ldots,n\})\,(i\neq j\Rightarrow a_{ij}=0)$.
\item The matrix $A$ is \hlt{upper triangular} iff $(\forall\,i,j\in\{1,\ldots,n\})\,(i>j\Rightarrow a_{ij}=0)$.
\end{enumerate}
\end{definition}

\begin{theorem} \label{thm:complex-operator-has-eigenvalue} \hlt{Complex Operators Have Eigenvalues}\\
Let $\mathcal{V}$ be a nonzero finite-dimensional complex vector space. Every $T\in\LL(\mathcal{V})$ has an eigenvalue.
\begin{proof}
Fix $T\in\LL(\mathcal{V})$. Choose a nonconstant annihilating polynomial for $T$ by Proposition \ref{prop:finite-dimensional-operators-have-annihilating-polynomials}.\\
One of its complex roots is an eigenvalue of $T$ by Theorem \ref{thm:annihilating-polynomials-control-spectrum}(ii).\\
\hlt{Structural Commentary}: An annihilator and complex factorization reduce existence to a singular linear factor.\\
\hlt{Key Insight}: A zero product of shifted operators contains a singular factor.
\end{proof}
\end{theorem}

\begin{theorem} \label{thm:upper-triangular-criteria} \hlt{Criteria for an Upper-Triangular Operator Matrix}\\
Let $n\in\mathbb{Z}_{>0}$, let $T\in\LL(\mathcal{V})$, and let $\beta=(\mathbf{v}_1,\ldots,\mathbf{v}_n)$ be a basis of $\mathcal{V}$. The following are equivalent:
\begin{enumerate}[label=(\roman*)]
\setlength{\itemsep}{0pt}
\item $[T]_{\beta\leftarrow\beta}$ is upper triangular;
\item $(\forall\,j\in\{1,\ldots,n\})\,T\mathbf{v}_j\in\operatorname{span}(\mathbf{v}_1,\ldots,\mathbf{v}_j)$;
\item $(\forall\,j\in\{1,\ldots,n\})\,\operatorname{span}(\mathbf{v}_1,\ldots,\mathbf{v}_j)\in\operatorname{Inv}(T)$.
\end{enumerate}
\begin{proof}
\begin{enumerate}[label=(\roman*)]
\setlength{\itemsep}{0pt}
\item[$(i)$] $\Rightarrow (ii)$ Fix $j\in\{1,\ldots,n\}$. The $j$th column of $[T]_{\beta\leftarrow\beta}$ is $[T\mathbf{v}_j]_{\beta}$.\\
Its entries below row $j$ vanish by (i), so $T\mathbf{v}_j\in\operatorname{span}(\mathbf{v}_1,\ldots,\mathbf{v}_j)$. Thus (ii) holds by the arbitrariness of $j$.
\item[$(ii)$] $\Rightarrow (iii)$ Fix $j\in\{1,\ldots,n\}$ and $\mathbf{u}\in\operatorname{span}(\mathbf{v}_1,\ldots,\mathbf{v}_j)$.\\
Choose $a_1,\ldots,a_j\in\mathbb{F}$ with $\mathbf{u}=\sum\limits_{k=1}^{j}a_k\mathbf{v}_k$. Then $T\mathbf{u}=\sum\limits_{k=1}^{j}a_kT\mathbf{v}_k\in\operatorname{span}(\mathbf{v}_1,\ldots,\mathbf{v}_j)$ by (ii).\\
Hence the span is invariant under $T$ by Definition \ref{def:invariant-subspace}, and (iii) follows.
\item[$(iii)$] $\Rightarrow (i)$ For $j\in\{1,\ldots,n\}$, the vector $\mathbf{v}_j$ belongs to the $T$-invariant space $\operatorname{span}(\mathbf{v}_1,\ldots,\mathbf{v}_j)$ by (iii).\\
Thus $T\mathbf{v}_j\in\operatorname{span}(\mathbf{v}_1,\ldots,\mathbf{v}_j)$ by Definition \ref{def:invariant-subspace}. Every matrix entry below the diagonal therefore vanishes, so (i) holds.
\end{enumerate}
\hlt{Structural Commentary}: The implication cycle translates column support into preservation of every initial span.\\
\hlt{Key Insight}: Upper triangularity is the coordinate form of an invariant flag.
\end{proof}
\end{theorem}

\begin{corollary} \label{cor:upper-triangular-orthonormal-basis} \hlt{Upper-Triangular Form in an Orthonormal Basis}\\
Let $\mathcal{V}$ be a nonzero finite-dimensional inner-product space and let $T\in\LL(\mathcal{V})$.\\
If $T$ has an upper-triangular matrix relative to a basis of $\mathcal{V}$, then it has one relative to an orthonormal basis.
\begin{proof}
Choose a basis $\beta=(\mathbf{v}_1,\ldots,\mathbf{v}_n)$ for which $[T]_{\beta\leftarrow\beta}$ is upper triangular, and set $\mathcal{U}_j\defeq\operatorname{span}(\mathbf{v}_1,\ldots,\mathbf{v}_j)$.\\
Then $(\forall\,j\in\{1,\ldots,n\})\,\mathcal{U}_j\in\operatorname{Inv}(T)$ by Theorem \ref{thm:upper-triangular-criteria}(i),(iii). Let $\gamma=(\mathbf{e}_1,\ldots,\mathbf{e}_n)$ be the Gram--Schmidt output for $\beta$.\\
List $\gamma$ is orthonormal by Theorem \ref{thm:gram-schmidt-process}(ii). Prefix identities $(\forall\,j\in\{1,\ldots,n\})\,\operatorname{span}(\mathbf{e}_1,\ldots,\mathbf{e}_j)=\mathcal{U}_j$ hold by Theorem \ref{thm:gram-schmidt-process}(i).\\
Since $\operatorname{span}(\gamma)=\mathcal{U}_n=\mathcal{V}$, $\gamma$ is an orthonormal basis. Its matrix for $T$ is upper triangular by Theorem \ref{thm:upper-triangular-criteria}(iii),(i).\\
\hlt{Structural Commentary}: Gram--Schmidt transports the invariant flag to an orthonormal basis.\\
\hlt{Key Insight}: Triangularity depends on initial spans, not their chosen generators.
\end{proof}
\end{corollary}

\begin{theorem} \label{thm:complex-upper-triangularization} \hlt{Complex Upper Triangularization}\\
Let $\mathcal{V}$ be a nonzero finite-dimensional complex vector space and let $T\in\LL(\mathcal{V})$.\\
Then $T$ has an upper-triangular matrix relative to a basis of $\mathcal{V}$.
\begin{proof}
\begin{enumerate}[label=(\roman*)]
\setlength{\itemsep}{0pt}
\item \hlt{Induction}: Induct on $n\defeq\dim_{\mathbb{C}}(\mathcal{V})$. For $n=1$, every operator matrix is upper triangular. Assume $n>1$ and that every operator on a nonzero complex vector space of dimension less than $n$ has an upper-triangular matrix relative to a basis.
\item \hlt{Eigenline and Quotient}: Choose eigenvector $\mathbf{v}_1$ of $T$ by Theorem \ref{thm:complex-operator-has-eigenvalue}, set $\mathcal{U}\defeq\operatorname{span}(\mathbf{v}_1)$; $\mathcal{U}\in\operatorname{Inv}(T)$ by Proposition \ref{prop:one-dimensional-invariant-subspaces-and-eigenvectors}.\\
The relations $T/\mathcal{U}\in\LL(\mathcal{V}/\mathcal{U})$ and $\dim_{\mathbb{C}}(\mathcal{V}/\mathcal{U})=n-1>0$ hold by Proposition \ref{prop:restriction-and-quotient-operators-well-defined}(ii) and Theorem \ref{thm:dimension-of-quotient}(ii).\\
By induction, choose a quotient basis $(\mathbf{v}_2+\mathcal{U},\ldots,\mathbf{v}_n+\mathcal{U})$ in which $T/\mathcal{U}$ is upper triangular.
\item \hlt{Lifted Basis}: Choose representatives $\mathbf{v}_2,\ldots,\mathbf{v}_n\in\mathcal{V}$.\\
A zero relation among $\mathbf{v}_1,\ldots,\mathbf{v}_n$ descends to the quotient basis, so the coefficients of $\mathbf{v}_2,\ldots,\mathbf{v}_n$ vanish.\\
The remaining coefficient vanishes because $\mathbf{v}_1\neq\mathbf{0}$. Thus $\beta\defeq(\mathbf{v}_1,\ldots,\mathbf{v}_n)$ is a basis by Proposition \ref{prop:full-dimension-basis-criteria}(ii).
\item \hlt{Triangularity}: The relation $T\mathbf{v}_1\in\operatorname{span}(\mathbf{v}_1)$ holds by $\mathcal{U}\in\operatorname{Inv}(T)$. Fix $j\in\{2,\ldots,n\}$.\\
Choose $c_{2,j},\ldots,c_{j,j}\in\mathbb{C}$ satisfying $(T/\mathcal{U})(\mathbf{v}_j+\mathcal{U})=\sum\limits_{k=2}^{j}c_{k,j}(\mathbf{v}_k+\mathcal{U})$ by Theorem \ref{thm:upper-triangular-criteria}(i),(ii).\\
The relation $T\mathbf{v}_j-\sum\limits_{k=2}^{j}c_{k,j}\mathbf{v}_k\in\mathcal{U}=\operatorname{span}(\mathbf{v}_1)$ holds by Definition \ref{def:restriction-and-quotient-operators}(ii).\\
Thus $(\forall\,j\in\{1,\ldots,n\})\,T\mathbf{v}_j\in\operatorname{span}(\mathbf{v}_1,\ldots,\mathbf{v}_j)$, and $[T]_{\beta\leftarrow\beta}$ is upper triangular by Theorem \ref{thm:upper-triangular-criteria}(ii),(i).
\end{enumerate}
\hlt{Structural Commentary}: Induction lifts a triangular quotient basis above a complex eigenline.\\
\hlt{Key Insight}: Each lifted quotient relation differs by a vector in the eigenline.
\end{proof}
\end{theorem}

\begin{definition} \label{def:complexification-real-vector-space} \hlt{Complexification of a Real Vector Space}\\
Let $\mathcal{V}$ be a real vector space.
\begin{enumerate}[label=(\roman*)]
\setlength{\itemsep}{0pt}
\item The \hlt{complexification} of $\mathcal{V}$ is $\mathcal{V}_{\mathbb{C}}\defeq\mathcal{V}\times\mathcal{V}$, with $(\mathbf{u},\mathbf{v})$ denoted $\mathbf{u}+i\mathbf{v}$.
\item Addition is $(\mathbf{u}_1+i\mathbf{v}_1)+(\mathbf{u}_2+i\mathbf{v}_2)\defeq(\mathbf{u}_1+\mathbf{u}_2)+i(\mathbf{v}_1+\mathbf{v}_2)$. For $a,b\in\mathbb{R}$, define $(a+bi)(\mathbf{u}+i\mathbf{v})\defeq(a\mathbf{u}-b\mathbf{v})+i(b\mathbf{u}+a\mathbf{v})$.
\item The \hlt{complexification map} is $\operatorname{cpx}:\mathcal{V}\to\mathcal{V}_{\mathbb{C}}$, $\operatorname{cpx}(\mathbf{u})\defeq\mathbf{u}+i\mathbf{0}$. Its image is $\operatorname{im}(\operatorname{cpx})=\{\mathbf{u}+i\mathbf{0}:\mathbf{u}\in\mathcal{V}\}$.
\item Define $\operatorname{Re},\operatorname{Im}:\mathcal{V}_{\mathbb{C}}\to\mathcal{V}$ by $\operatorname{Re}(\mathbf{u}+i\mathbf{v})\defeq\mathbf{u}$ and $\operatorname{Im}(\mathbf{u}+i\mathbf{v})\defeq\mathbf{v}$. Then $\mathbf{z}=\mathbf{0}$ iff $\operatorname{Re}\mathbf{z}=\operatorname{Im}\mathbf{z}=\mathbf{0}$.
\end{enumerate}
\end{definition}

\begin{proposition} \label{prop:complexification-vector-space-basis} \hlt{Complexification Preserves Bases}\\
Let $\mathcal{V}$ be a finite-dimensional real vector space and let $\beta=(\mathbf{v}_1,\ldots,\mathbf{v}_n)$ be a real basis of $\mathcal{V}$. Set $\beta_{\mathbb{C}}\defeq\operatorname{cpx}(\beta)$. Then:
\begin{enumerate}[label=(\roman*)]
\setlength{\itemsep}{0pt}
\item $\mathcal{V}_{\mathbb{C}}$ is a complex vector space;
\item $\beta_{\mathbb{C}}$ is a complex basis of $\mathcal{V}_{\mathbb{C}}$;
\item $\dim_{\mathbb{C}}(\mathcal{V}_{\mathbb{C}})=\dim_{\mathbb{R}}(\mathcal{V})$.
\end{enumerate}
\begin{proof}
\begin{enumerate}[label=(\roman*)]
\setlength{\itemsep}{0pt}
\item Addition makes $\mathcal{V}_{\mathbb{C}}$ an abelian group with zero $\mathbf{0}+i\mathbf{0}$ and inverse $-(\mathbf{u}+i\mathbf{v})=(-\mathbf{u})+i(-\mathbf{v})$ by the real vector-space axioms.\\
Fix $\alpha=a+bi$, $\eta=c+di\in\mathbb{C}$ and $\mathbf{z}=\mathbf{u}+i\mathbf{v}\in\mathcal{V}_{\mathbb{C}}$. The real parts of $\alpha(\eta\mathbf{z})$ and $(\alpha\eta)\mathbf{z}$ both equal $(ac-bd)\mathbf{u}-(ad+bc)\mathbf{v}$.\\
Imaginary parts both equal $(ad+bc)\mathbf{u}+(ac-bd)\mathbf{v}$, so scalar multiplication is associative. The identity $1\mathbf{z}=\mathbf{z}$ follows directly.\\
The two distributive laws hold by Definition \ref{def:complexification-real-vector-space}(ii) and real vector-space distributivity. Hence $\mathcal{V}_{\mathbb{C}}$ is a complex vector space.
\item For $\mathbf{u},\mathbf{v}\in\mathcal{V}$, choose real coefficients satisfying $\mathbf{u}=\sum\limits_{j=1}^{n}a_j\mathbf{v}_j$ and $\mathbf{v}=\sum\limits_{j=1}^{n}b_j\mathbf{v}_j$.\\
Then $\mathbf{u}+i\mathbf{v}=\sum\limits_{j=1}^{n}(a_j+ib_j)\operatorname{cpx}(\mathbf{v}_j)$, so $\beta_{\mathbb{C}}$ spans $\mathcal{V}_{\mathbb{C}}$. Suppose $\sum\limits_{j=1}^{n}(c_j+id_j)\operatorname{cpx}(\mathbf{v}_j)=\mathbf{0}$.\\
The equality $\sum\limits_{j=1}^{n}c_j\mathbf{v}_j=\sum\limits_{j=1}^{n}d_j\mathbf{v}_j=\mathbf{0}$ holds by Definition \ref{def:complexification-real-vector-space}(ii),(iv).\\
Thus $(\forall\,j\in\{1,\ldots,n\})\,c_j=d_j=0$ by independence of $\beta$, and $\beta_{\mathbb{C}}$ is a complex basis.
\item The basis $\beta_{\mathbb{C}}$ has $n=\dim_{\mathbb{R}}(\mathcal{V})$ members. Hence $\dim_{\mathbb{C}}(\mathcal{V}_{\mathbb{C}})=\dim_{\mathbb{R}}(\mathcal{V})$.
\end{enumerate}
\hlt{Structural Commentary}: Component identities establish scalar action; coefficient separation transports the basis and dimension.\\
\hlt{Key Insight}: Each complex coefficient is two independent real coefficients.
\end{proof}
\end{proposition}

\begin{corollary} \label{cor:complexification-coordinate-compatibility} \hlt{Coordinates Survive Complexification}\\
Let $\mathcal{V}$ be a finite-dimensional real vector space and let $\beta=(\mathbf{v}_1,\ldots,\mathbf{v}_n)$ be a basis of $\mathcal{V}$.\\
Set $\beta_{\mathbb{C}}\defeq\operatorname{cpx}(\beta)$. Then $(\forall\,\mathbf{v}\in\mathcal{V})\,[\operatorname{cpx}(\mathbf{v})]_{\beta_{\mathbb{C}}}=[\mathbf{v}]_{\beta}\in\mathbb{R}^n\subseteq\mathbb{C}^n$.
\begin{proof}
Fix $\mathbf{v}\in\mathcal{V}$ and write $[\mathbf{v}]_{\beta}=(a_1,\ldots,a_n)^t\in\mathbb{R}^n$. Then $\operatorname{cpx}(\mathbf{v})=\sum\limits_{j=1}^{n}a_j\operatorname{cpx}(\mathbf{v}_j)$ by Definition \ref{def:complexification-real-vector-space}(ii)--(iii).\\
Thus $[\operatorname{cpx}(\mathbf{v})]_{\beta_{\mathbb{C}}}=(a_1,\ldots,a_n)^t=[\mathbf{v}]_{\beta}$ by Proposition \ref{prop:complexification-vector-space-basis}(ii).\\
\hlt{Structural Commentary}: The embedding preserves one real expansion inside the complexified basis.\\
\hlt{Key Insight}: Embedded real vectors retain their real coordinate columns.
\end{proof}
\end{corollary}

\begin{proposition} \label{prop:complexification-real-dimension-and-embedding} \hlt{Complexification as a Real Space}\\
Let $\mathcal{V}$ be finite-dimensional real vector space, $\beta=(\mathbf{v}_1,\ldots,\mathbf{v}_n)$ be basis of $\mathcal{V}$. Regard $\mathcal{V}_{\mathbb{C}}$ as real vector space by restriction of scalars.\\
Set $\mathbf{e}_j\defeq\operatorname{cpx}(\mathbf{v}_j)$ $(j\in\{1,\ldots,n\})$. Define $\delta\defeq(\mathbf{e}_1,\ldots,\mathbf{e}_n,i\mathbf{e}_1,\ldots,i\mathbf{e}_n)$. Then:
\begin{enumerate}[label=(\roman*)]
\setlength{\itemsep}{0pt}
\item $\delta$ is a real basis of $\mathcal{V}_{\mathbb{C}}$;
\item $\dim_{\mathbb{R}}(\mathcal{V}_{\mathbb{C}})=2\dim_{\mathbb{R}}(\mathcal{V})$;
\item $\operatorname{cpx}:\mathcal{V}\to\mathcal{V}_{\mathbb{C}}$ is real-linear and injective.
\end{enumerate}
\begin{proof}
\begin{enumerate}[label=(\roman*)]
\setlength{\itemsep}{0pt}
\item For $\mathbf{u}+i\mathbf{w}\in\mathcal{V}_{\mathbb{C}}$, choose real coefficients with $\mathbf{u}=\sum\limits_{j=1}^{n}a_j\mathbf{v}_j$ and $\mathbf{w}=\sum\limits_{j=1}^{n}b_j\mathbf{v}_j$.\\
Then $\mathbf{u}+i\mathbf{w}=\sum\limits_{j=1}^{n}a_j\mathbf{e}_j+\sum\limits_{j=1}^{n}b_j(i\mathbf{e}_j)$, so $\delta$ spans over $\mathbb{R}$.\\
If this right side is zero, its real and imaginary parts vanish by Definition \ref{def:complexification-real-vector-space}(iv).\\
Thus $(\forall\,j\in\{1,\ldots,n\})\,a_j=b_j=0$ by independence of $\beta$, and $\delta$ is a real basis.
\item The basis $\delta$ has $2n=2\dim_{\mathbb{R}}(\mathcal{V})$ members.
\item For $r,s\in\mathbb{R}$ and $\mathbf{u},\mathbf{v}\in\mathcal{V}$, $\operatorname{cpx}(r\mathbf{u}+s\mathbf{v})=r\operatorname{cpx}(\mathbf{u})+s\operatorname{cpx}(\mathbf{v})$ by Definition \ref{def:complexification-real-vector-space}(ii)--(iii).\\
If $\operatorname{cpx}(\mathbf{u})=\operatorname{cpx}(\mathbf{v})$, then $\operatorname{cpx}(\mathbf{u}-\mathbf{v})=\mathbf{0}$ by linearity and $\mathbf{u}-\mathbf{v}=\mathbf{0}$ by Definition \ref{def:complexification-real-vector-space}(iv). Hence $\operatorname{cpx}$ is real-linear and injective.
\end{enumerate}
\hlt{Structural Commentary}: Real--imaginary separation builds the doubled basis and detects the embedding kernel.\\
\hlt{Key Insight}: Complexification is two independent real copies of $\mathcal{V}$.
\end{proof}
\end{proposition}

\begin{definition} \label{def:complexification-real-operator} \hlt{Complexification of a Real Operator}\\
Let $\mathcal{V}$ be a real vector space and let $T\in\LL(\mathcal{V})$. The \hlt{complexification} of $T$ is $T_{\mathbb{C}}:\mathcal{V}_{\mathbb{C}}\to\mathcal{V}_{\mathbb{C}}$, defined by $T_{\mathbb{C}}(\mathbf{u}+i\mathbf{v})\defeq T\mathbf{u}+iT\mathbf{v}$.
\end{definition}

\begin{proposition} \label{prop:complexification-operator-basic-properties} \hlt{Complexification Preserves Operator Data}\\
Let $\mathcal{V}$ be a real vector space and let $T\in\LL(\mathcal{V})$.
\begin{enumerate}[label=(\roman*)]
\setlength{\itemsep}{0pt}
\item $T_{\mathbb{C}}\in\LL(\mathcal{V}_{\mathbb{C}})$.
\item $(\forall\,k\in\mathbb{Z}_{\geq0})(\forall\,\mathbf{u},\mathbf{v}\in\mathcal{V})\,(T_{\mathbb{C}})^k(\mathbf{u}+i\mathbf{v})=T^k\mathbf{u}+iT^k\mathbf{v}$;
\item For a real basis $\beta=(\mathbf{v}_1,\ldots,\mathbf{v}_n)$ of $\mathcal{V}$, set $\beta_{\mathbb{C}}\defeq\operatorname{cpx}(\beta)$. Then $[T_{\mathbb{C}}]_{\beta_{\mathbb{C}}\leftarrow\beta_{\mathbb{C}}}=[T]_{\beta\leftarrow\beta}$ as complex matrices.
\end{enumerate}
\begin{proof}
\begin{enumerate}[label=(\roman*)]
\setlength{\itemsep}{0pt}
\item Let $\mathbf{z}=\mathbf{u}+i\mathbf{v}$, $\mathbf{w}=\mathbf{x}+i\mathbf{y}$, and $\alpha=a+bi\in\mathbb{C}$. Then $T_{\mathbb{C}}(\mathbf{z}+\mathbf{w})=T_{\mathbb{C}}\mathbf{z}+T_{\mathbb{C}}\mathbf{w}$ by the additivity of $T$.\\
By the real linearity of $T$, $T_{\mathbb{C}}(\alpha\mathbf{z})=T(a\mathbf{u}-b\mathbf{v})+iT(b\mathbf{u}+a\mathbf{v})=(aT\mathbf{u}-bT\mathbf{v})+i(bT\mathbf{u}+aT\mathbf{v})$.\\
The last term equals $\alpha T_{\mathbb{C}}\mathbf{z}$ by Definition \ref{def:complexification-real-vector-space}(ii). Hence $T_{\mathbb{C}}$ is complex-linear.
\item Fix $\mathbf{u},\mathbf{v}\in\mathcal{V}$. The identity holds for $k=0$. If it holds at $k$, then $(T_{\mathbb{C}})^{k+1}(\mathbf{u}+i\mathbf{v})=T_{\mathbb{C}}(T^k\mathbf{u}+iT^k\mathbf{v})=T^{k+1}\mathbf{u}+iT^{k+1}\mathbf{v}$.\\
Thus $(\forall\,k\in\mathbb{Z}_{\geq0})\,(T_{\mathbb{C}})^k(\mathbf{u}+i\mathbf{v})=T^k\mathbf{u}+iT^k\mathbf{v}$ by induction, and claim (ii) follows by the arbitrariness of $\mathbf{u},\mathbf{v}$.
\item The list $\beta_{\mathbb{C}}$ is a complex basis by Proposition \ref{prop:complexification-vector-space-basis}(ii). Fix $j\in\{1,\ldots,n\}$ and write $T\mathbf{v}_j=\sum\limits_{r=1}^{n}a_{r,j}\mathbf{v}_r$ with $a_{r,j}\in\mathbb{R}$.\\
Then $T_{\mathbb{C}}\operatorname{cpx}(\mathbf{v}_j)=\sum\limits_{r=1}^{n}a_{r,j}\operatorname{cpx}(\mathbf{v}_r)$ by Definition \ref{def:complexification-real-operator}. Hence these matrices have same columns by arbitrariness of $j$.
\end{enumerate}
\hlt{Structural Commentary}: Scalar expansion proves linearity, induction controls powers, and basis images identify columns.\\
\hlt{Key Insight}: Complexification changes the scalar field without changing real coordinate columns.
\end{proof}
\end{proposition}

Figure \ref{fig:complexification-real-operator} illustrates the componentwise action established in Proposition \ref{prop:complexification-operator-basic-properties}.\\
The proposition, rather than the diagram, owns complex linearity and coordinate compatibility.

\begin{illustrationfigure}[H]
  \centering
  \begin{tikzpicture}[
  >=Stealth,
  every node/.style={font=\small},
  box/.style={draw=gray!55, rounded corners=3pt, inner xsep=7pt, inner ysep=4.5pt, align=center},
  arrow/.style={->, line width=0.85pt, blue!70!black}
]
  \node[box] (real) at (0,0) {$\mathbf{u},\mathbf{v}\in\mathcal{V}$};
  \node[box] (complex) at (2.65,0) {$\mathbf{u}+i\mathbf{v}\in\mathcal{V}_{\mathbb{C}}$};
  \node[box] (image) at (5.80,0) {$T\mathbf{u}+iT\mathbf{v}$};

  \draw[arrow] (real) -- (complex);
  \draw[arrow] (complex) -- (image);
\end{tikzpicture}

  \caption{Complexification packages two real vectors: $T_{\mathbb{C}}(\mathbf{u}+i\mathbf{v})=T\mathbf{u}+iT\mathbf{v}$.}
  \label{fig:complexification-real-operator}
\end{illustrationfigure}

\begin{proposition} \label{prop:real-eigenvalues-complexification} \hlt{Real Eigenvalues of a Complexification}\\
Let $\mathcal{V}$ be a finite-dimensional real vector space, let $T\in\LL(\mathcal{V})$, and let $\lambda\in\mathbb{R}$. Then $\lambda$ is an eigenvalue of $T$ iff it is an eigenvalue of $T_{\mathbb{C}}$.
\begin{proof}
\begin{enumerate}[label=\roman*.]
\setlength{\itemsep}{0pt}
\item[$(\Rightarrow)$] Suppose $T\mathbf{v}=\lambda\mathbf{v}$ with $\mathbf{v}\neq\mathbf{0}$. Then $\operatorname{cpx}(\mathbf{v})\neq\mathbf{0}$ by Proposition \ref{prop:complexification-real-dimension-and-embedding}(iii).\\
The identity $T_{\mathbb{C}}\operatorname{cpx}(\mathbf{v})=\lambda\operatorname{cpx}(\mathbf{v})$ holds by Definition \ref{def:complexification-real-operator}.\\
Hence $\lambda$ is an eigenvalue of $T_{\mathbb{C}}$ by Definition \ref{def:eigenvalue-eigenvector}.
\item[$(\Leftarrow)$] Suppose $\mathbf{u},\mathbf{v}\in\mathcal{V}$ satisfy $T_{\mathbb{C}}(\mathbf{u}+i\mathbf{v})=\lambda(\mathbf{u}+i\mathbf{v})$ and $\mathbf{u}+i\mathbf{v}\neq\mathbf{0}$.\\
Then $T\mathbf{u}=\lambda\mathbf{u}$ and $T\mathbf{v}=\lambda\mathbf{v}$ by Definitions \ref{def:complexification-real-vector-space}(ii),(iv) and \ref{def:complexification-real-operator}.\\
The relation $\mathbf{u}\neq\mathbf{0}\lor\mathbf{v}\neq\mathbf{0}$ holds by Definition \ref{def:complexification-real-vector-space}(iv).\\
Thus one of $\mathbf{u},\mathbf{v}$ is an eigenvector of $T$ for $\lambda$ by Definition \ref{def:eigenvalue-eigenvector}.
\end{enumerate}
\hlt{Structural Commentary}: The embedding preserves a real eigenvector; component separation recovers one from its complexification.\\
\hlt{Key Insight}: A real eigenvalue acts identically on both real components.
\end{proof}
\end{proposition}

\begin{computationmethod} \label{meth:complexify-real-matrix} \hlt{Real Invariant Subspace from a Complex Eigenpair}\\
Given $n\in\mathbb{Z}_{>0}$ and $A\in\operatorname{Mat}(n,n,\mathbb{R})$, construct a nonzero subspace $\mathcal{U}\leq\mathbb{R}^n$ satisfying $A\mathcal{U}\subseteq\mathcal{U}$ and $\dim(\mathcal{U})\in\{1,2\}$:
\begin{enumerate}[label=(\roman*)]
\setlength{\itemsep}{0pt}
\item Let $T_A\in\LL(\mathbb{C}^n)$ be defined by $T_A\mathbf{z}\defeq A\mathbf{z}$.
\item Choose an eigenvalue $a+bi\in\mathbb{C}$ of $T_A$ by Theorem \ref{thm:complex-operator-has-eigenvalue}.\\
Choose an associated eigenvector $\mathbf{0}\neq\mathbf{u}+i\mathbf{v}\in\mathbb{C}^n$ by Definition \ref{def:eigenvalue-eigenvector}.\\
Then $A(\mathbf{u}+i\mathbf{v})=(a+bi)(\mathbf{u}+i\mathbf{v})$ by the definition of $T_A$.
\item Separate components to obtain $A\mathbf{u}=a\mathbf{u}-b\mathbf{v}$ and $A\mathbf{v}=b\mathbf{u}+a\mathbf{v}$.
\item Set $\mathcal{U}\defeq\operatorname{span}(\mathbf{u},\mathbf{v})$. Then $\mathcal{U}\neq\{\mathbf{0}\}$ by step (ii), and $A\mathcal{U}\subseteq\mathcal{U}$ by step (iii) and linearity.\\
A basis sublist of $(\mathbf{u},\mathbf{v})$ exists by Theorem \ref{thm:reducing-spanning-lists}, so $\dim(\mathcal{U})\in\{1,2\}$.
\end{enumerate}
\end{computationmethod}

\begin{theorem} \label{thm:real-operator-invariant-line-or-plane} \hlt{Invariant Line or Plane over $\mathbb{R}$}\\
Let $\mathcal{V}$ be a nonzero finite-dimensional real vector space and let $T\in\LL(\mathcal{V})$.\\
Then $(\exists\,\mathcal{U}\in\operatorname{Inv}(T))\,\dim(\mathcal{U})\in\{1,2\}$.
\begin{proof}
\begin{enumerate}[label=(\roman*)]
\setlength{\itemsep}{0pt}
\item \hlt{Complex Eigenpair}: The space $\mathcal{V}_{\mathbb{C}}$ is nonzero and finite-dimensional over $\mathbb{C}$ by Proposition \ref{prop:complexification-vector-space-basis}(i),(iii).\\
The complexified operator satisfies $T_{\mathbb{C}}\in\LL(\mathcal{V}_{\mathbb{C}})$ by Proposition \ref{prop:complexification-operator-basic-properties}(i). Choose an eigenvalue $\lambda\defeq a+bi\in\mathbb{C}$ by Theorem \ref{thm:complex-operator-has-eigenvalue}.\\
Choose $\mathbf{z}\defeq\mathbf{u}+i\mathbf{v}\in\mathcal{V}_{\mathbb{C}}\setminus\{\mathbf{0}\}$ satisfying $T_{\mathbb{C}}\mathbf{z}=\lambda\mathbf{z}$ by Definition \ref{def:eigenvalue-eigenvector}.
\item \hlt{Invariant Real Span}: The component identities $T\mathbf{u}=a\mathbf{u}-b\mathbf{v}$ and $T\mathbf{v}=b\mathbf{u}+a\mathbf{v}$ hold by Definitions \ref{def:complexification-real-vector-space}(ii),(iv) and \ref{def:complexification-real-operator}.\\
Set $\mathcal{U}\defeq\operatorname{span}(\mathbf{u},\mathbf{v})$. For $c,d\in\mathbb{R}$, $T(c\mathbf{u}+d\mathbf{v})=(ca+db)\mathbf{u}+(-cb+da)\mathbf{v}\in\mathcal{U}$ by the component identities.\\
Thus $\mathcal{U}\in\operatorname{Inv}(T)$ by Definition \ref{def:invariant-subspace}.
\item \hlt{Dimension}: The relation $\mathcal{U}\neq\{\mathbf{0}\}$ holds by $\mathbf{u}+i\mathbf{v}\neq\mathbf{0}$ and Definition \ref{def:complexification-real-vector-space}(iv).\\
A basis sublist of $(\mathbf{u},\mathbf{v})$ exists by Theorem \ref{thm:reducing-spanning-lists}. Consequently, $\dim(\mathcal{U})\in\{1,2\}$.
\end{enumerate}
\hlt{Structural Commentary}: A complex eigenpair separates into two generators whose span is closed under $T$.\\
\hlt{Key Insight}: The real and imaginary components form a nonzero invariant span.
\end{proof}
\end{theorem}

Figure \ref{fig:real-invariant-plane-from-complex-eigenpair} depicts the component separation in Theorem \ref{thm:real-operator-invariant-line-or-plane}.\\
The theorem, rather than the diagram, owns the invariant-subspace conclusion.

\begin{illustrationfigure}[H]
  \centering
  \begin{tikzpicture}[
  >=Stealth,
  every node/.style={font=\small},
  box/.style={draw=gray!55, rounded corners=3pt, inner xsep=7pt, inner ysep=4pt, align=center},
  arrow/.style={->, line width=0.85pt, blue!70!black}
]
  \node[box] (eig) at (0,0) {$T_{\mathbb{C}}(\mathbf{u}+i\mathbf{v})$\\$=(a+bi)(\mathbf{u}+i\mathbf{v})$};
  \node[box] (parts) at (3.55,0) {$T\mathbf{u}=a\mathbf{u}-b\mathbf{v}$\\$T\mathbf{v}=b\mathbf{u}+a\mathbf{v}$};
  \node[box] (plane) at (6.90,0) {$\mathcal{U}=\operatorname{span}(\mathbf{u},\mathbf{v})$\\$T\mathcal{U}\subseteq\mathcal{U}$};

  \draw[arrow] (eig) -- (parts);
  \draw[arrow] (parts) -- (plane);
\end{tikzpicture}

  \caption{A complex eigenpair determines a nonzero invariant real line or plane.}
  \label{fig:real-invariant-plane-from-complex-eigenpair}
\end{illustrationfigure}

\begin{definition} \label{def:almost-upper-triangular-real-matrix} \hlt{Almost Upper-Triangular Real Matrices}\\
Let $n\in\mathbb{Z}_{>0}$, $A\in\operatorname{Mat}(n,n,\mathbb{R})$. Matrix $A$ is \hlt{almost upper triangular} iff it admits a block representation with following properties:
\begin{enumerate}[label=(\roman*)]
\setlength{\itemsep}{0pt}
\item The block count and sizes satisfy $m\in\mathbb{Z}_{>0}$ and $(d_1,\ldots,d_m)\in\{1,2\}^m$.
\item One has $A=[A_{ij}]_{i,j=1}^{m}$, $A_{ij}\in\operatorname{Mat}(d_i,d_j,\mathbb{R})$, and $\sum\limits_{j=1}^{m}d_j=n$.
\item $(\forall\,i,j\in\{1,\ldots,m\})\,(i>j\Rightarrow A_{ij}=0)$.
\item Every diagonal block is either $A_{jj}=[\lambda_j]$ with $\lambda_j\in\mathbb{R}$, or $A_{jj}=\left[\begin{smallmatrix}a_j&-b_j\\b_j&a_j\end{smallmatrix}\right]$ with $a_j,b_j\in\mathbb{R}$ and $b_j\neq0$.
\end{enumerate}
The matrix $A$ is in \hlt{real Schur block form} iff it admits a block representation satisfying (i)--(iii), with no restriction beyond the prescribed sizes on its diagonal blocks.
\end{definition}

\begin{theorem} \label{thm:real-almost-upper-triangularization} \hlt{Real Almost Upper Triangularization}\\
Let $\mathcal{V}$ be nonzero finite-dimensional real vector space, $T\in\LL(\mathcal{V})$. Then $T$ has almost-upper-triangular matrix relative to basis of $\mathcal{V}$.
\begin{proof}
\begin{enumerate}[label=(\roman*)]
\setlength{\itemsep}{0pt}
\item \hlt{Induction}: Induct on $n\defeq\dim_{\mathbb{R}}(\mathcal{V})$. If $n=1$, every operator matrix is an admissible real $1\times1$ block. Assume $n\geq2$ and that every operator on a nonzero real vector space of dimension less than $n$ has an almost-upper-triangular matrix relative to a basis.
\item \hlt{Complex Eigenpair}: The space $\mathcal{V}_{\mathbb{C}}$ is a nonzero $n$-dimensional complex vector space by Proposition \ref{prop:complexification-vector-space-basis}(i),(iii).\\
The complexified operator satisfies $T_{\mathbb{C}}\in\LL(\mathcal{V}_{\mathbb{C}})$ by Proposition \ref{prop:complexification-operator-basic-properties}(i). Choose an eigenvalue $\lambda\defeq a+bi\in\mathbb{C}$ by Theorem \ref{thm:complex-operator-has-eigenvalue}.\\
Choose $\mathbf{z}\defeq\mathbf{u}+i\mathbf{v}\in\mathcal{V}_{\mathbb{C}}\setminus\{\mathbf{0}\}$ satisfying $T_{\mathbb{C}}\mathbf{z}=\lambda\mathbf{z}$ by Definition \ref{def:eigenvalue-eigenvector}.\\
Then $T\mathbf{u}=a\mathbf{u}-b\mathbf{v}$ and $T\mathbf{v}=b\mathbf{u}+a\mathbf{v}$ by Definitions \ref{def:complexification-real-vector-space}(ii),(iv) and \ref{def:complexification-real-operator}.
\item \hlt{Initial Invariant Block}: Construct $\mathcal{U}$ and its basis $\alpha$ by cases.
\begin{enumerate}[label=(\alph*)]
\setlength{\itemsep}{0pt}
\item If $b=0$, then $\mathbf{u}\neq\mathbf{0}\lor\mathbf{v}\neq\mathbf{0}$ by Definition \ref{def:complexification-real-vector-space}(iv). Choose $\mathbf{x}\in\{\mathbf{u},\mathbf{v}\}\setminus\{\mathbf{0}\}$. Set $\mathcal{U}\defeq\operatorname{span}(\mathbf{x})$ and $\alpha\defeq(\mathbf{x})$.\\
The component identities imply $T\mathbf{x}=a\mathbf{x}$, so $\mathcal{U}\in\operatorname{Inv}(T)$ and $[T|_{\mathcal{U}}]_{\alpha\leftarrow\alpha}=[a]$.
\item Suppose $b\neq0$. Implication $\mathbf{u}=\mathbf{0}\Rightarrow b\mathbf{v}=\mathbf{0}\Rightarrow\mathbf{v}=\mathbf{0}$ contradicts chosen eigenvector by first component identity.\\
Hence $\mathbf{u}\neq\mathbf{0}$. If $\mathbf{v}=c\mathbf{u}$ for $c\in\mathbb{R}$, then $c(a-bc)\mathbf{u}=cT\mathbf{u}=T\mathbf{v}=(b+ac)\mathbf{u}$ by linearity and the component identities.\\
Thus $b(c^2+1)=0$ by $\mathbf{u}\neq\mathbf{0}$, contrary to $b\neq0$. Therefore $\mathbf{u},\mathbf{v}$ are independent. Set $\mathcal{U}\defeq\operatorname{span}(\mathbf{u},\mathbf{v})$ and $\alpha\defeq(\mathbf{u},-\mathbf{v})$.\\
Then $\mathcal{U}\in\operatorname{Inv}(T)$ and $[T|_{\mathcal{U}}]_{\alpha\leftarrow\alpha}=\left[\begin{smallmatrix}a&-b\\b&a\end{smallmatrix}\right]$ by the component identities.
\end{enumerate}
In both cases, $\alpha$ is a basis of $\mathcal{U}$, $\dim(\mathcal{U})\in\{1,2\}$, and $[T|_{\mathcal{U}}]_{\alpha\leftarrow\alpha}$ is an admissible first diagonal block.
\item \hlt{Quotient Induction}: If $\mathcal{U}=\mathcal{V}$, the basis $\alpha$ proves the claim. Suppose $\mathcal{U}\neq\mathcal{V}$. Then $T/\mathcal{U}\in\LL(\mathcal{V}/\mathcal{U})$ by Proposition \ref{prop:restriction-and-quotient-operators-well-defined}(ii).\\
The inequality $0<\dim_{\mathbb{R}}(\mathcal{V}/\mathcal{U})=n-\dim_{\mathbb{R}}(\mathcal{U})<n$ holds by Theorem \ref{thm:dimension-of-quotient}(ii). Set $q\defeq n-\dim_{\mathbb{R}}(\mathcal{U})$.\\
By induction, choose a quotient basis $\overline{\gamma}\defeq(\mathbf{w}_1+\mathcal{U},\ldots,\mathbf{w}_q+\mathcal{U})$ in which $T/\mathcal{U}$ is almost upper triangular.
\item \hlt{Lifted Block Form}: Set $\gamma\defeq(\mathbf{w}_1,\ldots,\mathbf{w}_q)$ and $\beta\defeq(\alpha,\gamma)$.\\
A zero relation in $\beta$ descends to the quotient basis, so every coefficient of $\gamma$ vanishes.\\
The coefficients of $\alpha$ then vanish because $\alpha$ is a basis of $\mathcal{U}$. Thus $\beta$ is independent and has $n$ members.\\
Hence $\beta$ is a basis by Proposition \ref{prop:full-dimension-basis-criteria}(ii), and $\mathcal{V}=\mathcal{U}\oplus\operatorname{span}(\gamma)$.\\
The matrix of $T$ in $\beta$ is block upper triangular by Proposition \ref{prop:invariant-subspaces-and-block-form}(i).\\
Its first diagonal block is $[T|_{\mathcal{U}}]_{\alpha\leftarrow\alpha}$. Its lower-right block is $[T/\mathcal{U}]_{\overline{\gamma}\leftarrow\overline{\gamma}}$ by Definition \ref{def:restriction-and-quotient-operators}(i)--(ii).\\
Therefore $[T]_{\beta\leftarrow\beta}$ is almost upper triangular by Definition \ref{def:almost-upper-triangular-real-matrix}.
\end{enumerate}
\hlt{Structural Commentary}: Induction lifts an almost-triangular quotient above an invariant line or plane.\\
\hlt{Key Insight}: Complex eigenpairs supply the admissible real diagonal blocks.
\end{proof}
\end{theorem}

\begin{corollary} \label{cor:real-schur-block-form-orthonormal} \hlt{Orthonormal Real Schur Block Form}\\
Let $\mathcal{V}$ be nonzero finite-dimensional real inner-product space, $T\in\LL(\mathcal{V})$.\\
Then $T$ has real Schur block form relative to an orthonormal basis of $\mathcal{V}$.
\begin{proof}
\begin{enumerate}[label=(\roman*)]
\setlength{\itemsep}{0pt}
\item \hlt{Invariant Block Flag}: Choose a basis $\beta=(\mathbf{v}_1,\ldots,\mathbf{v}_n)$ in which $T$ is almost upper triangular by Theorem \ref{thm:real-almost-upper-triangularization}.\\
Fix its block sizes $(d_1,\ldots,d_m)\in\{1,2\}^m$. Set $s_k\defeq\sum\limits_{j=1}^{k}d_j$ and $\mathcal{U}_k\defeq\operatorname{span}(\mathbf{v}_1,\ldots,\mathbf{v}_{s_k})$.\\
The zero blocks below the block diagonal imply $(\forall\,k\in\{1,\ldots,m\})\,\mathcal{U}_k\in\operatorname{Inv}(T)$ by Definitions \ref{def:almost-upper-triangular-real-matrix}(iii) and \ref{def:invariant-subspace}.
\item \hlt{Orthonormal Flag}: Let $\gamma=(\mathbf{e}_1,\ldots,\mathbf{e}_n)$ be the Gram--Schmidt output for $\beta$. The list $\gamma$ is orthonormal by Theorem \ref{thm:gram-schmidt-process}(ii).\\
The prefix identities $(\forall\,k\in\{1,\ldots,m\})\,\operatorname{span}(\mathbf{e}_1,\ldots,\mathbf{e}_{s_k})=\mathcal{U}_k$ hold by Theorem \ref{thm:gram-schmidt-process}(i).\\
Since $\operatorname{span}(\gamma)=\mathcal{U}_m=\mathcal{V}$, $\gamma$ is an orthonormal basis.
\item \hlt{Block Form}: Relations $(\forall\,k\in\{1,\ldots,m\})\,\operatorname{span}(\mathbf{e}_1,\ldots,\mathbf{e}_{s_k})\in\operatorname{Inv}(T)$ make $[T]_{\gamma\leftarrow\gamma}$ block upper triangular for sizes $(d_1,\ldots,d_m)$.\\
Hence it is in real Schur block form by Definition \ref{def:almost-upper-triangular-real-matrix}.
\end{enumerate}
\hlt{Structural Commentary}: Gram--Schmidt preserves each subspace in the invariant block flag.\\
\hlt{Key Insight}: Orthonormalization changes block entries without changing block boundaries.
\end{proof}
\end{corollary}

\begin{proposition} \label{prop:upper-triangular-invertibility-criterion} \hlt{Invertibility from an Upper-Triangular Matrix}\\
Let $\dim(\mathcal{V})=n\in\mathbb{Z}_{>0}$ and $T\in\LL(\mathcal{V})$. Suppose $T$ has an upper-triangular matrix with diagonal $(\lambda_1,\ldots,\lambda_n)$ relative to a basis of $\mathcal{V}$.\\
Then $T$ is invertible iff $(\forall\,j\in\{1,\ldots,n\})\,\lambda_j\neq0$.
\begin{proof}
Let $\beta=(\mathbf{v}_1,\ldots,\mathbf{v}_n)$ be such a basis.
\begin{enumerate}[label=\roman*.]
\setlength{\itemsep}{0pt}
\item[$(\Rightarrow)$] Suppose $T$ is invertible and $(\exists\,j\in\{1,\ldots,n\})\,\lambda_j=0$; fix such a $j$. Set $\mathcal{U}_k\defeq\operatorname{span}(\mathbf{v}_1,\ldots,\mathbf{v}_k)$ and $\mathcal{U}_0\defeq\{\mathbf{0}\}$.\\
The first $j$ matrix columns have zero entries in rows $j,\ldots,n$ by upper triangularity and $\lambda_j=0$, so $T\mathcal{U}_j\subseteq\mathcal{U}_{j-1}$.\\
Map $S:\mathcal{U}_j\to\mathcal{U}_{j-1}$, $S\mathbf{u}\defeq T\mathbf{u}$, is therefore linear. Since $\dim(\mathcal{U}_j)=j>j-1=\dim(\mathcal{U}_{j-1})$, $S$ is not injective by Corollary \ref{cor:map-to-smaller-space-not-injective}.\\
Thus $(\exists\,\mathbf{u}\in\mathcal{U}_j\setminus\{\mathbf{0}\})\,T\mathbf{u}=\mathbf{0}$, contrary to the injectivity of $T$ by Theorem \ref{thm:finite-dimensional-operator-invertible-equivalences}(i),(ii). Hence $(\forall\,j\in\{1,\ldots,n\})\,\lambda_j\neq0$.
\item[$(\Leftarrow)$] Suppose $(\forall\,j\in\{1,\ldots,n\})\,\lambda_j\neq0$, and set $M\defeq[T]_{\beta\leftarrow\beta}$. Then $M$ is invertible by Theorem \ref{thm:inverse-upper-triangular}(i).\\
Choose $S\in\LL(\mathcal{V})$ satisfying $[S]_{\beta\leftarrow\beta}=M^{-1}$ by Theorem \ref{thm:linear-maps-and-matrices-same-data}(i).\\
The identity $[ST]_{\beta\leftarrow\beta}=M^{-1}M=I_n=[I_{\mathcal{V}}]_{\beta\leftarrow\beta}$ holds by Proposition \ref{prop:composition-and-matrix-multiplication}.\\
The identity $[TS]_{\beta\leftarrow\beta}=MM^{-1}=I_n=[I_{\mathcal{V}}]_{\beta\leftarrow\beta}$ holds by the same proposition.\\
Hence $ST=TS=I_{\mathcal{V}}$ by Theorem \ref{thm:linear-maps-and-matrices-same-data}(i), so $T$ is invertible by Definition \ref{def:invertible-linear-map-and-isomorphism}.
\end{enumerate}
\hlt{Structural Commentary}: A zero diagonal compresses an initial span; nonzero diagonals transport a matrix inverse.\\
\hlt{Key Insight}: Nonzero diagonal entries are exactly the invertible triangular layers.
\end{proof}
\end{proposition}

\begin{proposition} \label{prop:eigenvalues-of-upper-triangular-operator} \hlt{Eigenvalues from an Upper-Triangular Matrix}\\
Let $\dim(\mathcal{V})=n\in\mathbb{Z}_{>0}$ and $T\in\LL(\mathcal{V})$. Suppose $T$ has an upper-triangular matrix with diagonal $(\lambda_1,\ldots,\lambda_n)$ relative to a basis of $\mathcal{V}$.\\
Then $\operatorname{spec}(T)=\{\lambda_1,\ldots,\lambda_n\}$.
\begin{proof}
Fix $\lambda\in\mathbb{F}$. The matrix of $T-\lambda I_{\mathcal{V}}$ in the same basis is upper triangular with diagonal $(\lambda_1-\lambda,\ldots,\lambda_n-\lambda)$.\\
Thus $T-\lambda I_{\mathcal{V}}$ is not invertible iff $(\exists\,j\in\{1,\ldots,n\})\,\lambda=\lambda_j$ by Proposition \ref{prop:upper-triangular-invertibility-criterion}.\\
The left condition is equivalent to $\lambda\in\operatorname{spec}(T)$ by Theorem \ref{thm:equivalent-eigenvalue-conditions}(i),(iv) and Definition \ref{def:spectrum-and-geometric-multiplicity}(i).\\
Hence $\operatorname{spec}(T)=\{\lambda_1,\ldots,\lambda_n\}$ by the arbitrariness of $\lambda$.\\
\hlt{Structural Commentary}: Scalar shifting reduces the spectrum calculation to triangular invertibility.\\
\hlt{Key Insight}: A shifted diagonal entry vanishes exactly at an eigenvalue.
\end{proof}
\end{proposition}

\begin{perspective} \label{per:upper-triangular-viewpoint} \hlt{Why Triangular Form Matters}\\
Theorems \ref{thm:complex-upper-triangularization} and \ref{thm:real-almost-upper-triangularization} encode an operator by a basis-adapted flag.\\
Proposition \ref{prop:eigenvalues-of-upper-triangular-operator} reads the spectrum from diagonal entries whenever triangular form exists.\\
Upper-block coupling records interactions that diagonalization would remove.\\
This viewpoint organizes spectral data without asserting diagonalizability.
\end{perspective}
